\documentclass[11pt]{article}
\usepackage[a4paper,margin=2.4cm]{geometry}
\usepackage{fontspec}
\usepackage{xeCJK}
\setmainfont{Noto Serif}
\setCJKmainfont{Microsoft JhengHei}
\setCJKsansfont{Microsoft JhengHei}
\setlength{\parindent}{2em}
\setlength{\parskip}{0.5em}
\setlength{\emergencystretch}{2em}
\begin{document}

\section*{28. 群特徵標與理想論}

\begin{center}
J. Ber. d. DMV 34 (1925), S. 144
\end{center}

\textbf{3. E. Noether，哥廷根：群特徵標與理想論。}

Frobenius 的群特徵標理論——也就是有限群表示的理論——被理解為一個完全可約環（現代通常稱半單純環），即群環，的理想論。因此，一般地發展了完全可約環在加法分解下的同構定理；把它表示為直接不可分解的單純理想的直和時，這種表示在同構意義下唯一，並且屬於同一個雙邊理想的不可分解單邊理想彼此同構。因此，若把同構的理想合併為一類，則不同的不可分解理想類的數目，正好等於雙邊不可分解理想的數目。

專門化到超複量的完全可約系統——特別是群環——時，不可分解理想類和不可約表示類彼此一一對應。進一步，雙邊不可分解理想和特徵標相互對應；而一個雙邊理想的加法不可分解分量的數目，等於這些不可分解理想的秩。這樣，Frobenius 理論就被納入了這一框架。

詳述將發表於 \emph{Math. Ann.}

\end{document}
